\begin{entry}{Q023}{The Wigner function of a continuous-variable hyperedge}

\entrymeta{
  topic   = {Continuous-variable hypergraph states; Wigner function; nullifiers; non-Gaussianity},
  est     = {$\sim$35 minutes --- labelled stretch},
  origin  = {Original. The states, the nullifiers and Lemma~1 are taken from the first
    source; neither source evaluates the Wigner function.},
  source  = {D.~W.~Moore, \emph{Quantum hypergraph states in continuous variables},
    Phys.~Rev.~A \textbf{100}, 062301 (2019), Eq.~(2) and Lemma~1; A.~Ravikumar,
    D.~W.~Moore and R.~Filip, \emph{Nonclassical nullifiers for quantum hypergraph
    states}, Quantum (2026), arXiv:2502.02226},
  url     = https://arxiv.org/abs/1909.03871,
  prereq  = {The single-mode Wigner function of \qref{Q022}; the Fourier representation
    of the delta function},
  status  = drafted,
}

\begin{problem}
\emph{Context.} Continuous-variable graph states are built by applying two-mode
interactions $e^{i\hat q_j\hat q_k}$ to momentum eigenstates. They are Gaussian, so they
cannot by themselves give a computational advantage. Hypergraph states add interactions
among three or more modes, such as $e^{i\hat q_1\hat q_2\hat q_3}$, and a lattice of such
hyperedges is a resource for universal computation using only homodyne measurements.
These states are characterised by their \emph{nullifiers}, operators that annihilate the
state, and experiments certify them through the variances of the nullifiers, which are
second moments. Neither the definition nor the certification refers to the Wigner
function. This exercise computes it for a single hyperedge and asks what it reveals.

\emph{Setting.} Take $k$ bosonic modes with $[\hat q_j,\hat p_l]=i\delta_{jl}$ and
$\hbar=1$. For a real polynomial $f$ on $\RR^{k}$ consider the (unnormalisable) state
\[
  \ket{\psi_f}:=e^{if(\hat q)}\ket{p=0}^{\otimes k},
  \qquad
  \braket{q}{\psi_f}=e^{if(q)} ,
\]
and extend the Wigner function of \qref{Q022} to $k$ modes,
\[
  W(q,p):=\frac{1}{\pi^{k}}\int_{\RR^{k}}d^{k}y\;e^{2ip\cdot y}\,
  \psi_f(q-y)\,\psi_f^{*}(q+y),
  \qquad q,p\in\RR^{k}.
\]
Since the state is not normalisable, $W$ is a distribution. You may manipulate the
integrals formally, exchanging their order and using
$\int_{\RR}dy\,e^{iky}=2\pi\,\delta(k)$. For $f=q_1q_2$ (one graph-state edge) and
$f=q_1q_2q_3$ (one hyperedge) the nullifiers are
\[
  \hat n_i:=\hat p_i-\prod_{j\neq i}\hat q_j ,
\]
and we write $n_i(q,p):=p_i-\prod_{j\neq i}q_j$ for the corresponding functions on phase
space. You may use the following identity, which follows from the integral
representations of the Bessel functions $Y_0$ and $K_0$ (NIST DLMF, \S10.9 and
\S10.32) after the substitution $t=e^{u}$. For real $a,b$ with $ab\neq0$,
\[
  \int_{0}^{\infty}\frac{dt}{t}\,\cos\!\Bigl(at+\frac{b}{t}\Bigr)
  =\begin{cases}
     -\pi\,Y_0\bigl(2\sqrt{ab}\bigr), & ab>0,\\[2pt]
     2\,K_0\bigl(2\sqrt{-ab}\bigr), & ab<0 .
   \end{cases}
\]
Here $K_0(x)>0$ for all $x>0$, decays like $e^{-x}$, and diverges like $-\ln x$ as
$x\to0^{+}$. $Y_0(x)$ diverges like $\tfrac{2}{\pi}\ln x$ as $x\to0^{+}$ and oscillates
with decaying amplitude; its first zeros are $x\approx0.894,\ 3.958,\ 7.086$, and it is
positive between the first two.

\begin{enumerate}[label=\textbf{(\alph*)},leftmargin=*]
  \item \emph{Graph states.} Show that
        \[
          W(q,p)=\frac{1}{\pi^{k}}\int d^{k}y\;
          \exp\!\Bigl(i\bigl[2p\cdot y-f(q+y)+f(q-y)\bigr]\Bigr).
        \]
        Evaluate $W$ for $f=q_1q_2$ and describe it in terms of the nullifiers.
        Explain in one sentence why the same structure holds for every quadratic $f$,
        and why such a $W$ is non-negative.

  \item \emph{The hyperedge.} Let $f=q_1q_2q_3$. Reduce $W$ to a one-dimensional
        integral, and show that $W$ depends on the six phase-space coordinates only
        through a single function of $n_1,n_2,n_3$.

  \item \emph{Where the non-Gaussianity lives.} Using the identity above, write $W$ in
        closed form. Identify the set on which $W$ is singular and the region where
        it is negative, both in terms of the nullifiers. Compare with (a): where was
        $W$ supported for the graph state, and where is it supported now? For contrast,
        the single-mode cubic phase state $e^{i\gamma\hat q^{3}}\ket{p=0}$ has a Wigner
        function that is an Airy function of $p-3\gamma q^{2}$. Which structural
        feature of $f=q_1q_2q_3$ leads to a different special function?

  \item \emph{Lemma~1 in phase space.} Measure $\hat q_3$ and obtain the outcome $m$.
        You may use that $\int dp_3\,W(q_1,q_2,m;p_1,p_2,p_3)$ is the Wigner function
        of the (unnormalised) conditional state of modes $1$ and $2$. Compute it, show
        that it is non-negative, and identify the state. Which term of the exponent in
        (a) carries the negativity of the three-mode $W$, and why does the integration
        over $p_3$ remove it?
\end{enumerate}
\end{problem}

\begin{solution}[Part (a)]
\end{solution}

\begin{solution}[Part (b)]
\end{solution}

\begin{solution}[Part (c)]
\end{solution}

\begin{solution}[Part (d)]
\end{solution}

\begin{takeaway}
\end{takeaway}

% \attempt{
%   date       = ,
%   time       = ,
%   difficulty = ,
%   howitwent  = ,
%   insight    = ,
%   note       = ,
% }

\end{entry}